\ifdefined\XeTeXrevision
  \IfFileExists{TS1cmr.fd}{\documentclass[11pt]{article}}
    {\documentclass[11pt,UTF8,scheme=plain,fontset=fandol]{ctexart}}
\else
  \documentclass[11pt]{article}
\fi
\usepackage{amsmath,amssymb,amsthm,booktabs,hyperref}
\usepackage[margin=1in]{geometry}
\newtheorem{theorem}{Theorem}
\newtheorem{lemma}[theorem]{Lemma}
\newtheorem{proposition}[theorem]{Proposition}
\DeclareMathOperator{\lcm}{lcm}
\DeclareMathOperator{\diag}{diag}
\newcommand{\vp}{v_p}
\newcommand{\vg}{v_p^{\mathrm G}}
\newcommand{\Q}{\mathbb Q}
\newcommand{\Z}{\mathbb Z}
\newcommand{\R}{\mathbb R}
% Replaced from the final audited static certificate when packaging.
% BEGIN INSTANCE
\newcommand{\LL}{50}
\newcommand{\UU}{45}
\newcommand{\ZZ}{40,40,40,41,43}
\newcommand{\DD}{0,0,0,0,0,2,4,5,5,6,8,9,9,11}
\newcommand{\hh}{15}
\newcommand{\GateA}{50}
\newcommand{\TT}{5}
\newcommand{\JJ}{3}
\newcommand{\BB}{9}
\newcommand{\gateExp}{105}
\newcommand{\finalR}{-7.867370}
\newcommand{\finalC}{7.640135}
\newcommand{\finalNet}{-0.227235}
\newcommand{\cutoff}{9/200}
\newcommand{\phaseBins}{32}
% END INSTANCE
\title{One of the six values \(\zeta(7),\zeta(9),\ldots,\zeta(17)\) is irrational}
\author{[Author / GitHub ID]}
\date{24 November 2026}
\begin{document}
\maketitle
\begin{abstract}
We give a computer-assisted proof that at least one of the six real
Riemann zeta values \(\zeta(7),\zeta(9),\zeta(11),\zeta(13),\zeta(15),\zeta(17)\)
is irrational. The method uses a growing determinant of fourth derivatives
of a nested-pole kernel. The arithmetic estimate combines folded
jet capacities and a periodic calculation at small primes. Non-vanishing
on an infinite arithmetic progression of auxiliary primes follows from
a confluent Vandermonde matrix and fixed binomial blocks. The accompanying
static certificate and independent rational/Arb verifier specify and
verify the numerical inequalities used in the argument.
\end{abstract}

\begin{theorem}
At least one of \(\zeta(7),\zeta(9),\zeta(11),\zeta(13),\zeta(15),\zeta(17)\)
is irrational.
\end{theorem}

Earlier fixed-window results include Zudilin's theorem on
\(\zeta(5),\zeta(7),\zeta(9),\zeta(11)\) \cite{zudilin} and the
seven-value result from \(\zeta(7)\) to \(\zeta(19)\) of Liu \cite{liu}.
We use the growing-determinant strategy of \cite{liu}, with a different
kernel and certificates. The parameters of the present proof are
\[
 \begin{gathered}
 L=\LL,\quad U=\UU,\quad R=U+L,\quad h=\hh n,\\
 (z_a)=(\ZZ),\qquad (d_b)=(\DD).
 \end{gathered}
 \tag{1}
\]
Here \(L,U,z_a,d_b\) are fixed integers, and \(n\) is the asymptotic index.
The certificate uses \(\theta=h/(Ln)\). We reserve \(\mathcal H_n\) for
the determinant matrix and \(K=2L\theta\) for the scaled rank budget.

\section{Kernel and polynomial determinant}
Write \(s=\min_a z_a\), \(c_n=(U+R)n/2+1\), and
\[
 a_b=(U+d_b)n+1,\qquad b_b=(R-d_b)n+1,\qquad
 \ell_b=b_b-a_b.
\]
Define
\[
 R_n(t)=2(t+c_n)
 \prod_{a=1}^5
 \frac{\prod_{j=1}^{z_an}(t+j)(t+2c_n-j)}{((z_an)!)^2}
 \prod_{b=0}^{13}\frac{\ell_b!}{\prod_{j=a_b}^{b_b}(t+j)} .
 \tag{2}
\]
This rational function is odd under \(t\mapsto-2c_n-t\).
Its zero order at each of \(-sn,\ldots,-1\) is at least five. No
pole lies at those integers. The parameters satisfy
\(\sum_b(L-2d_b)-2\sum_a z_a>4L\theta\). Thus the product of
\(R_n(t)\) with any two basis polynomials below is \(O(t^{-2})\)
at infinity, in fact with considerably greater decay.

For an integer center put \(E_0=1\) and, for \(i\ge1\),
\[
 E_i(x)=\binom{x+i}{2i}+\binom{x+i-1}{2i}.
\]
These even polynomials have degree \(2i\), leading coefficient
\(2/(2i)!\), and take integer values on \(\Z\).
In the half-integral case one can instead use
\(E_i(x)=\binom{x+i-1/2}{2i}\), with leading coefficient \(1/(2i)!\)
and the same evenness on the half-integral lattice.
For \(0\le i,j<h\) let
\[
 (\mathcal H_n(\boldsymbol\zeta))_{ij}
 =\frac1{4!}\sum_{t=-sn}^{\infty}
   (R_n(t)E_i(t+c_n)E_j(t+c_n))^{(4)}
 \tag{3}
\]
where the sum runs over integers. If \(t=-k+w\) is at a pole,
write the partial fractions of the product in (3) as
\(\sum_{\ell}a_{k,\ell}w^{-\ell}\) plus a regular part.
Differentiation and summation give a contribution
\[
 \sum_\ell \binom{\ell+3}{4}a_{k,\ell}
       \bigl(\zeta(\ell+4)-H_{\ell+4}(k-sn-1)\bigr),
 \quad H_j(M)=\sum_{v=1}^{M}v^{-j}.
 \tag{4}
\]
Reflection sends the coefficient at \(k\) to
\((-1)^{\ell+1}a_{k,\ell}\) at \(2c_n-k\). Thus even zeta orders
cancel. The sum of all \(a_{k,1}\) is zero by decay at infinity,
so the order-five zeta term also vanishes. The order-fourteen
poles give only the six claimed odd orders. Replacing these values
by independent indeterminates \(\mathbf X\) gives a matrix with
affine entries. Let
\(\det\mathcal H_n(\mathbf X)=\gamma_n P_n(\mathbf X)\), where
\(P_n\) is a primitive integer polynomial and \(\gamma_n\) is its
rational content (the sign can be fixed arbitrarily). For the zero
polynomial set \(P_n=0\) and \(\gamma_n=1\). The degree of \(P_n\)
is at most \(h\).

\section{Arithmetic bounds}
For a polynomial with rational coefficients we write \(\vg\)
for the minimum of their \(p\)-adic valuations, with the value
\(+\infty\) for zero. We use \(\vg(FG)=\vg(F)+\vg(G)\).
The arithmetic cost is \(-\log|\gamma_n|\).

\begin{lemma}[Taylor estimates]
Let \(p\) be odd, \(p^2>2i\), and let \(x\) be on the appropriate
integer or half-integer lattice. Writing \(E_i^{[r]}\) for a Taylor
coefficient, one has
\[
 \vp(p^r E_i^{[r]}(x))\ge
          \max\{0,r-\lfloor 2i/p\rfloor\}.
 \tag{5}
\]
For every nonnegative integer \(q\), adding
\(\max(\lfloor 2i/p\rfloor-q,0)\) to the left side gives a bound
of at least \(\max(r-q,0)\).
\end{lemma}
\begin{proof}
For integer \(A\), Vandermonde's identity expresses
\(\binom{A+w}{d}\) as \(\sum_{j=0}^d\binom A{d-j}\binom wj\).
The coefficient of \(w^r\) in \(\binom wj\) is \(s(j,r)/j!\),
where \(s(j,r)\) is a signed Stirling number of the first kind.
This gives valuation at least \(-\lfloor d/p\rfloor\) when
\(p^2>d\). Also, for \(j\ge1\),
\(\binom wj=(-1)^{j-1}(w/j)\prod_{k=1}^{j-1}(1-w/k)\).
Its coefficient of \(w^r\) is a sum of products with \(r\)
denominator factors, all at most \(j<p^2\), giving valuation at
least \(-r\). Apply these two bounds to the integer binomial
arguments in \(E_i\). For the last assertion use
\(\max(0,r-d)+\max(d-q,0)\ge\max(0,r-q)\).
\end{proof}

\begin{lemma}[All prime powers and the very small prime range]
There is a fixed \(M_0=O(n)\) such that
\(D_{M_0}^{18}\mathcal H_n(\mathbf X)\) has integer coefficients,
where \(D_m=\lcm(1,\ldots,m)\). Consequently the upper
contribution to \(-\log|\gamma_n|/(nhL)\) from primes
\(p\le\varepsilon n\) is at most \(18\varepsilon/L+o(1)\).
\end{lemma}
\begin{proof}
Take \(M_0\) to dominate all block lengths, all nonzero distances
from the poles to the factors, all harmonic cutoffs, and \(2h\).
The numerator blocks divided by their length factorials are
integer binomial polynomials at integral pole centers. The
factorization used above shows that a Taylor coefficient of
order \(r\) of any such binomial, or of the basis, has denominator
dividing \(D_{M_0}^r\), with no restriction on \(p\).

Fix \(P=p^e\). In a denominator block containing \(k\), the
number of remaining indices congruent to \(k\bmod P\) is
\(\lfloor(k-a_b)/P\rfloor+\lfloor(b_b-k)/P\rfloor\le
\lfloor\ell_b/P\rfloor\). In a block outside \(k\), the count of
congruent indices is at most \(\lfloor\ell_b/P\rfloor+1\).
The terms of size \(\ell_b/P\) therefore cancel against the
factorial normalization at every \(e\). For a numerator block
the count in its consecutive interval differs from its length
divided by \(P\) by at most one, and binomial integrality
provides the nonnegative constant-coefficient valuation after
its factorial normalization. This accounts explicitly for the
cancellation of all terms of order \(n/P\).

If \(m\) denominator blocks contain \(k\), the constant
coefficient of the local inverse denominator has denominator
dividing \(D_{M_0}^{14-m}\). Each additional Taylor order of
its nonvanishing factors costs at most one more copy of
\(D_{M_0}\), from expansion of \((u+w)^{-1}\) for \(u\ne 0\).
The odd central factor is an integral polynomial and does not
worsen the estimate. Hence the Laurent coefficient of order
\(\ell\) in (4) costs at most \(D_{M_0}^{14-\ell}\).
Multiplying by the harmonic term costs \(D_{M_0}^{\ell+4}\).
Leibniz expansion of the determinant proves
\(-\vg(\gamma_n)\le 18h\lfloor\log_p M_0\rfloor\).
For \(p\le\sqrt{M_0}\) the contribution divided by \(nhL\) is
at most \(18\sqrt{M_0}\log(M_0)/(nL)=o(1)\). For the remaining
primes up to \(\varepsilon n\) the exponent is at most one,
and the prime number theorem gives the claimed bound.
\end{proof}

\subsection{Local slot costs and reflection}
We now work away from the very small prime range, so that
\(p^2>M_0\). Pair the poles at \(k\ge c_n\) with their reflections,
and let \(m\) be the order before any cancellation at the center.
For a site let \(v\) be the valuation of the constant coefficient
of the regularized kernel (including the factorial normalization).
For a noncentral site let \(N_d,N_z\) count nonzero factors in
the denominator and numerator whose distances to the site are
multiples of \(p\), respectively. After removing the \(m\)
zero denominator factors, the change of valuation of the
coefficient of order \(d\) in the local regular factor is at
least \(-d\) if there is a denominator alias. Without denominator
aliases it is at least \(-\min(d,N_z)\): each vanishing numerator
factor is linear. The factor \(2(t+c_n)\) can either be counted
among these numerator factors or estimated as an integral
polynomial. In the limiting count below the residue class on
which it vanishes modulo \(p\) costs only \(O(1)\) slots per prime
and ignoring that zero can only weaken the bound.

Let \(r\) be the Taylor order falling on the product
\(E_iE_j\). For each \(\ell\) with \(1\le\ell\le m-r\), the
order in the regular factor is \(d=m-\ell-r\).
For \(\ell\) odd the paired harmonic coefficient can have
negative valuation from a multiple of \(p\) in either cutoff.
For \(\ell\) even, only the difference of the two harmonics
occurs. If their counts of multiples of \(p\) are equal, those
summands cancel exactly. Let \(I_\ell\) be the resulting
indicator for possible denominators. Therefore the
coefficient of \(E_i^{[a]}E_j^{[b]}\), \(a+b=r\), in the
paired form has Gauss valuation bounded below by the minimum
of \(v-d'-(\ell+4)I_\ell\), with \(d'=d\) or
\(\min(d,N_z)\) as above. After extracting the vectors
\(p^a E_i^{[a]}\) and \(p^b E_j^{[b]}\), a valid cost is
\[
 C_r=\max\bigl(0,\max_{1\le\ell\le m-r}
         (r-v+d'+(\ell+4)I_\ell)\bigr),
 \qquad B_a=\max_{r\ge a}C_r .
 \tag{6}
\]
Both \(B_a\) and \(B_b\) dominate \(C_{a+b}\).
The unrefined cost is \(m-v+4\) if the right cutoff contains
a multiple of \(p\), and \(m-v-1\) otherwise, cut off below by zero.

\subsection{Folded Cauchy--Binet}
Let \(Y_{i,(k,a)}=p^a E_i^{[a]}(c_n-k)\).
The preceding construction expresses the matrix as \(Y W Y^T\)
with block-diagonal \(W\). In each block an entry indexed by
\(a,b\) has Gauss valuation at least \(-(B_a+B_b)/2\).
Equivalently one can factor these weights in a ramified extension,
normalizing \(v_p(p)=1\). The coefficient inequalities over the
rational ground field are unchanged.

Group centers modulo \(p\) up to sign. If \(x\) is congruent to
a chosen representative \(z\) on the appropriate lattice, then
\[
 p^a E_i^{[a]}(x)=\sum_{j\ge a}\binom ja
       \bigl((x-z)/p\bigr)^{j-a}p^jE_i^{[j]}(z).
 \tag{7}
\]
Changing \(x\) to \(-x\) supplies only the sign \((-1)^a\).
Multiply row \(i\) of \(Y\) by
\(p^{(\lfloor2i/p\rfloor-q)_+}\). By (5) the degree-\(j\)
vector in (7) now has valuation at least \((j-q)_+\).
Any nonzero determinant of selected columns must choose
distinct such degrees within each residue class. Thus taking
\(s'\) columns of one class costs at most the sum of its
largest \(s'\) weights minus \(2\sum_{j<s'}(j-q)_+\) in the
two-sided determinant bound. More explicitly, each term of
the two Cauchy--Binet expansions uses two sets of \(h\)
columns. The valuation bound is the average of the costs of
the two sets together with the common row penalty. Taking the
largest possible cost on each side proves the same generous
bound for the Gauss valuation of the determinant.
For a class whose weights in decreasing order are \(B_j\),
the incremental costs are
\(B_j-2(j-q)_+\), indexed from \(j=0\).
Include the row penalty \(2\sum_{i<h}(\lfloor2i/p\rfloor-q)_+\).
One may also use the unfolded bound directly from (5), or
repeat the maximal \(B_j\) in a class to obtain a simpler folded
bound. For any set of nonnegative incremental costs the sum of
the largest \(h\) equals the sum, over positive integer levels,
of the number of increments at least that level capped at \(h\).

\subsection{Arithmetic integrals and stability at finite endpoints}
Use the coordinates \(x=2p/n\), \(u=2k/n\), and put \(N_0=2L\)
and \(K=N_0\theta\). The prime coordinate stored in the
certificate is \(x/N_0=p/(nL)\).
The limiting pole intervals are
\([2(U+d_b),2(R-d_b)]\), and the zero intervals are
\([0,2z_a]\), \([2(U+R-z_a),2(U+R)]\).
Their factorial lengths are the interval lengths in these
coordinates. At a given \(u\) counts of congruent sites in an
interval \([a,b]\) are
\(\lfloor(b-u)/x\rfloor-\lfloor(a-u)/x\rfloor\) away from
endpoints. The two harmonic cutoffs are
\(\lfloor(u-2s)/x\rfloor\) and
\(\lfloor(2(U+R)-u-2s)/x\rfloor\).
Subtract the factorial counts with the appropriate signs to
compute \(v\) in (6). On the folded axis
\(z=2\operatorname{dist}(u-(U+R),x\Z)\), the residue-class
interval is \(0\le z\le x\). A width on this axis is divided
by two to obtain residue-class capacity in the \(u\) scale.
At each level cap the capacity at \(K\). The continuous row
penalty is \(\sum_{j=q+1}^{\lfloor 2K/x\rfloor}(2K-jx)\).
Thus integration in \(x\) followed by division by \(N_0K\)
gives the contribution to the normalized arithmetic budget.

For a fixed cutoff \(p/n\ge\varepsilon>0\), only finitely
many lines in \(u\) and their folded images are relevant.
Their coefficients are fixed integers. The exact count
discontinuities and all intersections are affine.
The verifier checks ordering on each open rational slab by
signs at both endpoints of the slab. Equality at an endpoint
is permitted. Factorial walls and changes in \(\lfloor2K/x\rfloor\)
are also checked. For any certified \(q\) all direct
capacities are affine in a slab. The min operation at \(K\)
is integrated by splitting at the exact rational saturation
point. Different bounds can be chosen on subdivisions of a
slab. This construction justifies the chosen bounds without
requiring that the optimizing \(q\) be constant on a geometric
slab. The separate audit also computes the lower envelope by
solving the affine crossing equations and checks three interior
rational points per geometric slab.

For the profile (1) we have
\(R-s\ge\max(L,\max_a z_a,2h/n)\). If \(p>(R-s)n\),
the harmonic factors, pole-to-pole differences, and the Taylor
coefficients of the factorial-normalized binomial polynomials
are \(p\)-integral. The cost bound vanishes beyond this endpoint.

Here is the prelimit justification. For \(x\ge 2\varepsilon\)
the number of translated endpoints and possible folds is
bounded in terms of the fixed profile and \(\varepsilon\).
In the finite lattice each endpoint, center, or harmonic cut
is displaced in the \(u\) scale by \(O(1/n)\). Exclude primes
whose \(x\) lies within \(O(1/n)\) of a fixed factorial wall
or affine intersection. There are \(O_\varepsilon(1)\) such
integer prime candidates. The previous denominator lemma bounds
their total normalized effect by \(O_\varepsilon(\log n/n)\).
At every remaining prime, the displaced cutoff functions agree
with the limiting functions outside \(O_\varepsilon(1)\)
intervals of length \(O(1/n)\). Each interval contains
\(O_\varepsilon(1)\) lattice points. After folding, a residue
profile can change only at the images of these points or at the
finitely many corresponding shifted residue endpoints. Each
class has \(O_\varepsilon(1)\) sites, and each slot weight is
bounded by \(18\). Likewise an interval of folded width \(w\)
has \(nw/4+O(1)\) residue classes. Thus every level capacity
and its capped value have error \(O_\varepsilon(1)\) in the
unscaled lattice. The finite row penalty is the sum of a
bounded number of terms of the form
\((h-\lceil jp/2\rceil)_+\), and has the same error.
The resulting integrated error after division by \(nhL\)
is \(o(1)\) by the prime number theorem.
This argument gives an asymptotic bound, and allows a positive
finite endpoint error at an individual prime.
Finally the denominator lemma permits \(\varepsilon\downarrow0\).

\subsection{Periodic small prime certificate}
For small \(x\) we use a bound valid for all \(x\) in the indicated
small-prime interval. The unrefined weight at any site is
at most \(B=4+\#\text{pole factors}-\#\text{zero factors}-
\text{factorial valuation}\). Use \(B_+=\max(B,0)\).
Write \(1/x=k+\tau\) with \(0\le\tau<1\), and use residue
coordinate \(w\) modulo one, measured relative to the center.
The count from a signed interval with integer endpoints \(a,b\)
relative to that center is
\[
  \pm\big(\lfloor b\tau-w\rfloor-\lfloor a\tau-w\rfloor
                         -\lfloor(b-a)\tau\rfloor\big).
 \tag{8}
\]
The integer parts involving \(k\) cancel. By reflection the
whole residue-circle average equals the folded average.
Let \(F(\tau)=\int_0^1 B_+\,dw\) and
\(G(\tau)=\int_0^1\sum_{v=0}^{B_+-1}\lfloor v/2\rfloor\,dw\).
The affine wall equations in \(w,\tau\) certify piecewise
affine formulas for both quantities.

Write \(\eta=\{2K/x\}\), \(m=\lfloor2K/x\rfloor\).
For any integer \(0\le u_0\le m\) take \(q=m-u_0\).
Giving each residue class unlimited slots with the same coarse
weight and dropping the caps yields, per row, the upper bound
\[
 F(\tau)+\frac{x}{2K}
   \left(G(\tau)+(1-u_0)F(\tau)+u_0(u_0-1)
                          +\eta(2u_0-F(\tau))\right).
 \tag{9}
\]
Indeed the layer \(j\) has at most
\(q+1+\lfloor(B_+-j)/2\rfloor\) slots in each class, and one
then adds the row penalty. Since \(2K\) is integral for (1),
\(\eta=\{2K\tau\}\). The verifier checks each choice of \(u_0\)
against the maximum allowed by the cutoff and checks its quadratic
correction in (9). Substituting \(x=(k+\tau)^{-1}\), the
integrals of the affine and quadratic pieces over the first
45 periods use rational functions and logarithms. For later
periods take \(u_0=0\) and bound the correction above by \(F+G\).
If \(\overline F,\overline G\) are period integrals the tail is
bounded by \(\overline F/44+(\overline F+\overline G)/(4K\cdot44^2)\)
before division by \(N_0\). This uses integral comparison for
the inverse-square and inverse-cube series.

\section{Analytic certificate}
Let \(L_n=Ln\) and take the contour
\(\sigma_0=-\lfloor c L_n\rfloor-1/2\) for the rational \(c\)
in the certificate. The residues of
\((\pi\cot\pi t)^{(4)}(R_n P)(t)/4!\) at integers are
\((R_nP)^{(4)}(t)/4!\).
The rational kernel has no poles to the right of the contour.
The omitted integers between \(-sn\) and the contour give zero
because of the fivefold zeros. Closing the contour to the right
on half-integer vertical sides gives (3), up to an orientation
sign which is immaterial for the determinant modulus. The
horizontal edges vanish by exponential decay of the fourth
cotangent derivative. The right vertical edge vanishes by
rational decay and the integrable exponential bound in the
imaginary direction. For large positive imaginary part, the expression
\(\cot(\pi t)=i(e^{2\pi it}+1)/(e^{2\pi it}-1)\)
has a convergent geometric expansion in \(e^{2\pi it}\).
The fourth derivative kills the constant and is bounded by a
constant times \(e^{-2\pi\Im t}\). Conjugation gives the
bound in the lower half-plane. On a half-integer vertical line
the remaining bounded imaginary range stays a fixed distance
from the poles, uniformly by periodicity.
Thus the contour weight has magnitude
bounded by a fixed constant times \(|R_n(t)|e^{-2\pi|\Im t|}\).

For \(y=|\Im t|/L_n\) and \(\delta=(U+R)/(2L)-c\), the scalar
phase \(\phi_c(y)\) is obtained by summing
\[
 \int_a^b\log|u-c+iy|\,du-(b-a)\log(b-a)+(b-a)
 \tag{10}
\]
over the normalized zero intervals with sign \(+1\) and pole
intervals with sign \(-1\), and subtracting \(2\pi y\).
Stirling's estimate gives the normalization term. The
half-integer displacement gives a uniform \(O(\log L_n)\)
error in \(\log|R_n|\) on \(0\le y\le4\).
For completeness, on each side of \(u=c\) the logarithm is
monotone. Upper and lower Riemann sums outside the one or two
mesh cells nearest \(c\) differ only by endpoint terms of size
\(O(\log L_n)\) after multiplication by \(L_n\).
The exceptional cells have integral magnitude
\(O(\log L_n/L_n)\) and the half-integer contour keeps their
sampled logarithms above \(-O(\log L_n)\).
Shifting the center by \(O(1/L_n)\) or comparing the two
adjacent endpoint cells leaves the same bound.
The pole intervals stay a positive distance to the right of
\(c\). Formula (10) is evaluated via
\(X(d,y)=\frac d2\log(d^2+y^2)-y\arctan(y/d)\) for \(d>0\),
using the signed endpoints and the corresponding linear
\(\pi y/2\) contributions. The linear endpoint term from the
logarithmic antiderivative cancels the \(+(b-a)\) in (10).

\begin{lemma}[Discrete root polynomials]
Suppose a row band has positive rational weights \(w_a\)
summing to one, with real nonnegative atoms \(a\).
Write \(\delta_n=(c_n+\sigma_0)/L_n\) and
\[
 B_{a,n}(z)=\big(z-(\delta_n+ia)^2\big)\big(z-(\delta_n-ia)^2\big),
 \qquad
 Q_i(z)=z^{\,i-2\sum_a k_a}\prod_a B_{a,n}(z)^{k_a},
 \quad k_a=\lfloor iw_a/2\rfloor .
 \tag{11}
\]
The row norm at exponential scale is bounded above by the
continuous phase
\[
 \phi_c(y)+\frac{t}{2}\sum_a w_a\log M_a(y),
 \quad
 M_a(y)=(y^2-a^2)^2((y+a)^2+4\delta^2)((y-a)^2+4\delta^2),
 \tag{12}
\]
with a displacement \(t=i/L_n+O(1/L_n)\) and uniform
\(O(\log L_n/L_n)\) error in the upper comparison.
\end{lemma}
\begin{proof}
On the contour \(z=(\delta_n+iy)^2\) one has
\(|B_{a,n}(z)|^2=M_a(y)\) with \(\delta_n\) in place of
\(\delta\). The two strictly positive quadratic factors change
by a factor \(1+O(1/L_n)\) uniformly. The vanishing factor is
unchanged. Let \(M\) be a common denominator of the weights
in the band. With \(j=\lfloor i/(2M)\rfloor\) the integer
\(jMw_a\) is at most \(k_a\), and the excess factors and the
fixed-point remainder in (11) have total degree bounded in
terms of \(M\) alone. On \(0\le y\le4\) their modulus is
bounded uniformly from above. Use \(t=2Mj/L_n\) in (12).
This proves the stated one-sided comparison even at the
zeros. A two-sided unregularized comparison is unavailable
at a root with a changed multiplicity. Off neighborhoods
where \(M_a(y)<L_n^{-A}\), the floor-rounded expression
itself differs by \(O(\log L_n/L_n)\) for any fixed \(A>0\).
For the upper estimate over the entire contour the
whole-block comparison just proved is stronger. Only a
bounded number of rows per band can cross an endpoint because
of the displacement in \(t\). A uniform coarse bound of
\(O(L_n)\) per exceptional row follows from bounded atoms
and the same decreasing tail, so they do not affect the limit.
\end{proof}

The logarithmic moment bound for a row is at most the supremum
of (12). On any band the supremum of affine functions of
\(t\) is convex, so the integral over each band is bounded by
the trapezoid of the two certified endpoint suprema.
The verifier bounds the function on \([0,4]\) with Arb
intervals, using endpoint bounds for factors which vanish.
For \(y\ge4\), numerator derivative contributions are bounded
by \(11/y\), denominators have a nonpositive derivative
contribution, and the derivatives of the root terms are
bounded by \(5\theta/y\) (since \(2\sum_a z_a / L = 408/50 < 11\)
and \(4/(1 - (7/20)^2) < 5\), as all \(a\le7/5\)). Thus the derivative
of the upper tail phase is at most
\(-2\pi+(11+5\theta)/y<0\), bounded away from zero.
The same tail domination applies to the finite products,
with \(O(1/L_n)\) adjustments to the exponents, so the
tail integral satisfies the moment bound as well.

For the monic even basis \(Q_i((t+c_n)^2/L_n^2)\), the
Heine--Andr\'eief identity expresses the determinant as an
integral of a squared evaluation determinant against the
complex contour weight. Taking absolute values bounds it by
the corresponding integral with the squared modulus and
the positive modulus of the weight. By the same identity
this is a positive Gram determinant. Hadamard bounds it by
the product of the diagonal moment integrals.
Changing back to \(E_i\) multiplies the bound by
\(\prod_i(\operatorname{lc}(E_i)L_n^{2i})^2\).
Its log divided by \(nhL\) tends to
\(\theta(3-2\log(2\theta))\) by Stirling's formula. The
static certificate combines this term with the
\(\phaseBins\) verified phase bands.

\section{The confluent auxiliary-prime gate}
The following parameters are certified:
\(T=\TT\), \(J=\JJ\), \(A=\GateA\), \(B=\BB\),
where \(A=R-s-T\), \(h=JTn\), \(s+T\le U\),
\(s+2T\le L\), and \(B>2(J-1)\).
For \(p=An+1\) prime the rightmost \(Tn\) poles have exactly one
multiple of \(p\) in their right harmonic cutoff. There is no
other singular cutoff: \(p>Ln\), the left mirrored pole
is outside this tail, and the largest cutoff is less than \(2p\).
The factorials and pole-to-pole differences are units.
At these tail sites \(k-p\) runs through \(sn+1,\ldots,(s+T)n\).
The only numerator roots among nonzero multiples of \(p\)
at these sites thus come from the left blocks. The inequality
\(s+2T\le L\) rules out a multiple from the right blocks:
the maximum right-root distance is \((U+T)n-1<p\).
The central factor is a unit at the tail sites.

Let \(q\) be the left-root count at a tail site, and \(m\)
the local pole order. The chosen offsets match the jumps at
\(d=s+T-z_a\) for \(s<z_a<s+T\). The permanent hits \(z_a\ge s+T\)
are accounted for at offset zero. These equalities, verified
in the certificate, give \(m-q=B-4\).
The regular local numerator can therefore be written
\((w-p)^q U_k(w)\) with \(U_k(0)\) a unit. In (4) use the
geometric series \((1-w/p)^{-5}\) to sum the singular
summands \(p^{-\ell-4}\binom{\ell+3}{4}\).
For Taylor orders \(a,b\) the entry after multiplication by
\(p^{B-a-b}\) is congruent to
\[
 (-1)^{q+1}U_k(0)
       \binom{m+3-q-a-b}{m-1-a-b}\pmod p,
 \tag{13}
\]
with zero when \(a+b\ge m\). Terms coming from higher
coefficients of \(U_k\) acquire extra factors of \(p\).
The regular harmonics and zeta coordinates are integral.
The \(J\)-by-\(J\) binomial blocks in (13) have nonzero
integer determinants, explicitly listed in the certificate.

Form \(V\) whose columns are the first \(J\) Taylor vectors at
each of the \(Tn\) tail sites. It is square of size \(h\) and
invertible modulo \(p\). Indeed the squares of the \(x_k=c_n-k\)
are distinct and nonzero modulo \(p\), since \(2|x_k|<p\)
throughout the tail and the absolute values are distinct.
The change \(x\mapsto x^2\) is locally invertible at these
points. Thus the evaluation in Taylor jets of even polynomials
is the confluent interpolation isomorphism on polynomials of
degree \(<h\) in \(x^2\). The basis leading coefficients are
units since \(p>2(h-1)\).

Conjugate \(\mathcal H_n\) by \(V^{-1},V^{-T}\), and scale its
entry in retained jet indices \(i,j\in\{0,\ldots,J-1\}\) by
\(p^{B-i-j}\).
If the two contributing jet orders are both \(<J\), their
vectors are coordinate vectors after conjugation, and only
the same-site blocks from (13) remain. If \(a\ge J\) but \(b<J\),
the second vector fixes \(j=b\). The first yields at least
\(a-i\ge1\) extra powers of \(p\). The case \(a<J,b\ge J\) is
symmetric with \(b-j\ge1\). If both are at least \(J\),
their excess is at least two. Regular matrix coefficients
remain divisible by \(p\) after scaling, since \(B-i-j>0\).
These arguments are valid for every specialization of the
zeta indeterminates integral at \(p\). Except for the finitely
many primes dividing the finitely many binomial determinants,
the reduced matrix is invertible. In particular
\[
 v_p(\det\mathcal H_n(\mathbf X))=-\gateExp n
 \tag{14}
\]
at all such specializations. Dirichlet's theorem supplies
infinitely many primes \(p\equiv1\bmod A\).

\section{Certificate and conclusion}
The arithmetic and analytic estimates have the common
normalization \(nhL\). With the cutoff \(\cutoff\) in the
prime coordinate \(p/(nL)\), the certificate proves the
two upper bounds
\[
 \limsup\frac{\log|\det\mathcal H_n(\boldsymbol\zeta)|}{nhL}
       < \finalR,\qquad
 \limsup\frac{-\log|\gamma_n|}{nhL}<\finalC .
\]
The combined upper bound is at most \(\finalNet<0\), with
the displayed decimal bounds rounded outward.
If all six zeta values were rational, choose their common
denominator \(Q\). The numbers \(Q^hP_n(\boldsymbol\zeta)\)
would be integers. By (14) they would be nonzero for infinitely
many \(n\), after discarding primes dividing \(Q\). The
negative bound and \(h=O(n)\) make their absolute values
tend to zero. This is impossible.

The static data contain each phase cap, root atom and rational
weight, the middle/high-prime slab partition and selected
branches, the periodic small-prime tiles, and the local gate
blocks. Run \texttt{python verify\_zeta6\_certificate.py}
with \texttt{python-flint}. This verifier uses only rational
and Arb calculations. It does not need the optimizer. For
audit and Lean planning see the separate reproduction notes.

\begin{thebibliography}{9}
\bibitem{liu} J. Liu, \emph{At least one of \(\zeta(7),\ldots,\zeta(19)\)
is irrational}, 2026. \url{https://github.com/ALexxxxxLiu/odd-zeta-seven-lean}.
\bibitem{zudilin} W. Zudilin, \emph{One of the numbers \(\zeta(5),\zeta(7),
\zeta(9),\zeta(11)\) is irrational}, Russian Math. Surveys 56 (2001), 774--776.
\bibitem{flsz} S. Fischler, L. Lai, J. Sprang, W. Zudilin,
\emph{Many more odd zeta values are irrational}, arXiv:2611.16156 (2026).
\end{thebibliography}
\end{document}
